% ---
% title: Reward Hacking
% date: 2026/2/9
% tag: RL
% ---
\section{Flow-GRPO 里的 reward hacking}

用在线 RL（GRPO / PPO）训练时，模型经常通过取巧的方式提高 reward，而不是真正改善任务能力。

常见表现：贪婪地优化可测指标，图像质量反而崩坏；为了拿 reward 偏向某些风格、颜色、构图；在 preference reward 下只生成 evaluator 喜欢的模板式构图；多样性丢失，构图单一，模式坍缩。

RL 会强推模型往 reward 最大的局部最优去挤。缺了 restraint（限制），模型就会忽略通用能力和多样性，变得最优化但极度偏执。这就是 reward hacking。

\section{KL Regularization 的作用：把 RL 限在“微调区间”}

KL 项（通常和 pretrained model 的分布比）把新旧策略的差异压小：

\[
\text{KL}(\pi_{\text{new}} \,\|\, \pi_{\text{old}})
\]

它的含义是模型的输出分布不能离初始模型太远，更新要围绕原模型的小邻域进行。这有两个关键效果。

\subsection{保持 image quality}

原模型（如 SD3.5）在自然图像质量、构图、纹理上已经训练得很充分。没有 KL，RL 可以让 \(v_θ\) 大幅偏离原模型，很容易破坏生成能力，DrawBench 这类图像质量指标会下降。

\subsection{保持 model diversity}

KL 也限制新模型把概率质量全压到某些 seed、某些 style 上，或者把输出收缩到某个窄 distribution。

没有 KL 时，RL 很容易模式坍缩：不同 seed 生成近乎一样的图，因为 reward 最高的模式被强化了无数次。论文里说去掉 KL 后 “different seeds produce nearly identical results”，这就是典型的 RL-based mode collapse。

\section{不同任务里，去掉 KL 的坏法不一样}

组合生成和文本渲染这两类任务里，去掉 KL 的后果是图像质量明显下降，preference score 下降，还出现 reward hacking。

这些任务的 reward 和图像整体结构关系不大，大部分是局部判断：“两个红色苹果”“一只猫在桌子左边”、字体是否可读。低质量但高相关度的图就能骗过它们：更简单的构图、更少的背景、更强的对比度、不自然的渲染。所以去掉 KL，模型会训得很糟。

\section{PickScore 任务为什么看起来没事}

因为 PickScore 与最终评估指标高度一致。PickScore reward 本身就在让图像更美观、更清晰、更符合审美。

去掉 KL 后，模型不需要作弊，它真的能优化到 evaluator（PickScore）的偏好，图像质量变得更像 PickScore 喜欢的风格。

但问题还在。

\section{即使在 PickScore 里，样式也会坍缩}

PickScore 会偏好某种构图、某种色调、某种主体位置。去掉 KL，RL 不受约束，模型会倾向于生成 PickScore 最爱的那一种风格；即使 prompt 和 seed 不同，调性也相似。

多样性崩溃，图像看起来过度一致，模型退化成 PickScore 模板生成器。论文里说 “Outputs converge to a single style across different seeds.” 这说明 RL 把模型的分布压扁了，把所有概率质量挤到 reward 最高的局部区域。典型的 RL mode collapse。

\section{KL 在约束什么}

KL 不阻止模型更新。它做的是 constrained optimization：在不破坏已有能力的情况下往奖励方向推进。

不做 KL 就是 unconstrained optimization；加了 KL 就是 trust-region / small-step optimization。这也是 PPO/GRPO 背后最经典的思想之一。

Flow-GRPO 的任务 reward 决定往哪走，图像质量却来自训练前的基座模型。KL 维持的是基础能力和任务能力之间的平衡。
